\section{Front wheeler on a transformer}
\label{sec:exp3}

We now change both the architecture and the optimizer: from SGD-trained
MLPs to a transformer under AdamW.

\paragraph{Conditions.} Character-level language modeling on tiny
Shakespeare: a 10.8M-parameter decoder-only transformer (6 layers,
$d{=}384$, 6 heads, context 256, dropout $0.2$), AdamW ($3{\times}10^{-4}$,
200-step warmup, cosine decay, weight decay $0.1$, clip $1.0$), 5{,}000
steps, batch 64, 3 seeds. The front is the last transformer block, the
final LayerNorm, and the LM head; the LM head is \emph{untied} from the
token embedding so front steps cannot leak trunk updates through weight
tying. Front steps run the trunk under \texttt{no\_grad}. \fw{} keeps the
MLP configuration unchanged. Metric: validation loss (lower is better);
cost: backward compute as a fraction of \textsc{full}'s. The complete
population is Table~\ref{tab:exp3matrix} (Appendix~\ref{app:tables}).

\paragraph{The low-budget result transfers.} The claim is
that \fw{}'s matched-budget advantage is not an SGD/MLP artifact.
\textsc{full} reaches $1.488 \pm 0.008$ and \textsc{front} $1.836 \pm
0.004$; \fw{} at a $6\%$ deep-step budget ($22\%$ backward compute) reaches
$1.773 \pm 0.013$---indistinguishable within seed spread from \lpft{}
spending $50\%$ of steps ($1.766 \pm 0.007$; an $8\times$ ratio in
deep-step budget, ${\sim}2.7\times$ in backward compute)---robust across
thresholds ($1.773$--$1.776$ for $\tau$ varied $10\times$), with \bcd{}
tracking at $1.797 \pm 0.010$. At matched low compute, no fixed schedule
approaches this: \lpft{}($0.95$) reaches $1.835$ and \textsc{periodic}-20
$1.910$. The gap to always-full remains substantial: \fw{} recovers
${\sim}18\%$ of the front-to-full gap here, vs.\ ${\sim}$two-thirds on
MLPs---deep trunks need more deep training on harder tasks, and \fw{}
trades quality for compute rather than matching full BP.

\begin{figure}[t]
\centering
\includegraphics[width=0.72\textwidth]{exp3_pareto}
\caption{Transformer: validation loss vs.\ full-backprop budget (log
scale; mean $\pm$ std, 3 seeds; the \fw{} diamond pools thresholds,
whiskers span min to max). Periodic interleaving is
\emph{destructive}---worse than front-only at every budget, monotonically
in interleaving frequency---and one-switch \lpft{} tracks the front-only
line until half the budget is spent; at low budget, \fw{} sits below the
front-only line where one-switch and periodic schedules do not reach.}
\label{fig:exp3pareto}
\end{figure}

\paragraph{The governor is the budget dial.} Sweeping the trigger
threshold $\tau$ across two orders of magnitude ($0.005 \to 1.0$) with the
governor active barely moves the operating point: the backoff cancels eager
alarms and the budget stays pinned at $6$--$9\%$ of steps---the threshold
is a floor, not a dial. The dial is the governor itself: sweeping
$\kappa_{\max}$ from $2000$ down to $50$ traces a smooth, monotone Pareto
curve---budgets $6\% \to 49\%$ of steps, validation loss $1.773 \to
1.596$ (Figure~\ref{fig:exp3pareto})---that dominates \lpft{} and
\textsc{periodic}-$N$ at every tested budget on the transformer: at an
$8\%$ budget the dial already matches \lpft{} spending $50\%$. One
parameter selects the compute--quality operating point with no retuning of
the trigger. (On the CIFAR-10 MLP the dial's high-budget endpoint, a sweep
not reported in Section~\ref{sec:exp2}, is comparable to but does not beat
the best fixed $50\%$ schedules: $53.8$ vs.\ $54.2$ for \lpft{}.)

\paragraph{The freezing ratchet is strong but fragile.} At
\fw{}'s exact backward-compute budget, \textsc{ratchet}($0.1$) collapses to
$2.271 \pm 0.04$---worse than periodic---because it freezes five of six
blocks before the trunk has learned anything, and a ratchet cannot
un-freeze. \textsc{ratchet}($0.3$) does beat \fw{}'s default ($1.696$ vs.\
$1.773$) at $1.46\times$ the compute, and remains slightly ahead of the
\fw{} dial at matched mid-budget compute ($1.696$ at $32\%$ vs.\ $1.717$ at
$33.5\%$): a \emph{well-timed} ratchet is genuinely competitive in the
mid-budget regime. The pattern across testbeds is consistent: the ratchet's
timing is task-critical and unrecoverable when wrong, whereas \fw{} ran
identical trigger hyperparameters on every testbed in this paper.

\paragraph{Interleaving toxicity worsens under Adam.} It worsens. Periodic interleaving is destructive at every budget
tested---$1.87$ at $2\%$, worsening monotonically with interleaving
frequency to $2.40 \pm 0.05$ at a $20\%$ budget, far worse than doing
nothing deep---and it happens \emph{despite} gradient clipping, which cured
the SGD case: Adam's update $m/\sqrt{v}$ is approximately scale-invariant
in the gradient, so clipping the gradient does not clip the step. Damage
scales with interleaving \emph{frequency}, implicating repeated
perturbation---each isolated deep step moves trunk features under a head
sharply co-adapted to the old ones---and trunk moment estimates go stale
across gaps ($\beta_2 = 0.999$ means a 50-step burst refreshes only
${\sim}5\%$ of $v$), the schedule-induced analogue of the sparse-update
pathology that motivates SparseAdam for embedding tables.

\paragraph{Episodic optimizer state.}
Table~\ref{tab:optstate} is the decisive comparison, and it refined our
hypothesis. \textsc{reset-warmup} \emph{rescues} periodic interleaving
almost entirely ($2.05 \to 1.82$, from far worse than front-only to
slightly better): the dominant pathology of isolated deep steps under Adam
is stale-state re-entry at full aggression. But \textsc{reset} alone
\emph{backfires} ($2.05 \to 2.33$): a freshly initialized Adam step is the
most aggressive kind, so hygiene without gentleness is worse than
staleness. The surviving principle: \emph{the optimizer must fire with the
BP, and it must fire gently}---bursts amortize a warmup naturally, while
length-1 episodes can only approximate one by running deep steps at reduced
learning rate (for periodic, our 10-step warmup on length-1 episodes
reduces to exactly that, a caveat for interpretation). For \fw{} itself the
variants are a wash ($1.773$ persist vs.\ $1.773$ reset-warmup)---bursts
already re-equilibrate internally---which yields the systems result for
free: \emph{trunk optimizer state can be discarded between bursts at no
quality cost}, so Adam's two fp32 moments per trunk parameter need exist
only during the ${\sim}6\%$ of steps spent in bursts. The strict
\textsc{episodic} variant fails here ($2.13$), attributable to its SGD
front phases (transformers train poorly without adaptivity); a
stateless-but-adaptive front optimizer is the natural repair and is future
work.

\begin{table}[t]
\centering
\caption{Optimizer-state ablation (validation loss, mean $\pm$ std,
3 seeds). Anchors: front-only $1.836$, full BP $1.488$.}
\label{tab:optstate}
\begin{tabular}{lcc}
\toprule
Trunk Adam state & \fw{} ($6\%$) & \textsc{periodic}-10 ($10\%$) \\
\midrule
\textsc{persist} & $1.773 \pm 0.013$ & $2.047 \pm 0.127$ \\
\textsc{reset} & $1.782 \pm 0.013$ & $2.333 \pm 0.055$ \\
\textsc{reset-warmup} & $\mathbf{1.773 \pm 0.013}$ & $\mathbf{1.818 \pm 0.011}$ \\
\textsc{episodic} (SGD front) & $2.132 \pm 0.012$ & $2.437 \pm 0.008$ \\
\bottomrule
\end{tabular}
\end{table}

\subsection{Is the trigger necessary? Controls and distribution shifts}
\label{sec:controls}

Two simpler rivals could explain \fw{}'s results without any closed loop:
\textsc{ftlp} (spend all depth first, then never again) and \textsc{pburst}
(bursts on a fixed clock). We test both on the stationary task and under
two mid-training distribution shifts of different severity; complete
numbers are Tables~\ref{tab:exp3matrix} and~\ref{tab:shift}
(Appendix~\ref{app:tables}).

\paragraph{Deep-first and the fixed-cadence clock, stationary.} Only if
its one parameter is guessed correctly. At \fw{}'s matched budget,
\textsc{ftlp}($0.06$) is the worst cell in our study ($2.290 \pm 0.035$,
far below \textsc{front}): $300$ full steps from initialization---before
warmup completes---leave a half-trained trunk that is then frozen forever.
Given $3\times$ the budget, \textsc{ftlp}($0.2$) is the \emph{best} cheap
schedule we measured ($1.665 \pm 0.007$ at a $20\%$ deep-step budget).
Deep-first is a bet on knowing the right early budget: its peak and its
collapse are two settings of the same parameter. The fixed-cadence clock at
matched budget is consistently behind \fw{} ($1.814 \pm 0.014$ vs.\
$1.773 \pm 0.013$; $1.801 \pm 0.010$ with reset-warmup) though far from
failure; note its period was set \emph{from \fw{}'s own discovered
budget}---the clock is a hindsight-tuned baseline.

\paragraph{Two distribution shifts.} We switch the
training corpus mid-run at a step the schedules are not told, placed
mid-gap for \textsc{pburst}'s clock so it cannot be lucky: a \emph{mild}
shift (tiny Shakespeare $\to$ War and Peace) and a \emph{harsh} one ($\to$
C source code; the front-to-full gap grows from $0.35$ to $0.57$ of loss).
Under the mild shift \fw{} is the best cheap schedule ($1.549 \pm 0.013$),
ahead of the fair-placed clock in all three seed pairings ($1.561 \pm
0.008$), and its result is invariant to where the shift lands, which no
clock can claim; \textsc{ftlp} drops to the \textsc{front} floor ($1.605$
vs.\ $1.611$). Under the harsh shift, plain \fw{} ($1.449 \pm 0.005$)
again edges the fair-placed clock ($1.453$), while \textsc{ftlp} lands
between the floor and the burst schedules ($1.522 \pm 0.008$; its
prose-trained trunk transfers generic sequence features to code better than
a random trunk). The rivals' specific failure cells thus move with the
condition---which is itself the point: a fixed schedule's fate depends on
how well its frozen assumption happens to match the condition, and no
tested condition moved \fw{} out of the top of the cheap field or below
the \textsc{front} floor.

\paragraph{A shift-aware trigger.} \fw{}'s
documented weakness is that the backoff, stretched long by pre-shift
saturation, delays post-shift re-escalation (mild shift: $12\%$ of the
front-to-full gap recovered, vs.\ $18\%$ stationary). The repair is one
line: reset the cooldown when the loss EMA \emph{jumps above} its running
minimum---a shift signature, the opposite of the plateau the trigger
watches for. Under the harsh shift this \emph{governor reset} responds by
spending $6.7\%$ deep steps (vs.\ plain \fw{}'s $4\%$; the budget is
self-selected) and reaches $1.411 \pm 0.005$---the best cheap schedule
under the hardest condition we tested, recovering $24\%$ of the
front-to-full gap vs.\ plain \fw{}'s $18\%$. The trigger's verdict across
every test, stated at full strength and no more: bursts carry the
performance; a fixed clock, its period set with hindsight, ties or
slightly trails \fw{}; the closed loop's demonstrated products are
discovering the schedule without that hindsight, robustness when
assumptions break, and---with the one-line reset---responsiveness to them
breaking.

\paragraph{Measured resources validate the cost model---with one caveat.}
End-to-end single runs (2{,}000 steps, identical hardware;
Appendix~\ref{app:infra}) measure: full BP $98.4$s, \fw{} $49.3$s,
\textsc{ratchet}($0.1$) $49.7$s, front-only $46.6$s---\fw{} and the
compute-matched ratchet are equal-cost in wall clock to within $1\%$, and
\fw{} delivers a real $2.0\times$ end-to-end speedup over full BP at this
scale (per-step: $23.3$ms front vs.\ $49.2$ms full). The caveat is peak
activation memory: front-only peaks at $0.77$\,GB, but any schedule that
ever takes a full-depth step---\fw{} included---peaks at $3.03$\,GB,
because the peak is set by the deepest single step. \fw{} reduces
\emph{average} memory and (via episodic state) trunk optimizer memory, not
the activation peak; activation checkpointing inside bursts, which are
only ${\sim}6\%$ of steps, is the natural remedy and is future work.

\noindent The results so far take the premise---that cheap front steps
carry most per-step progress---as given; the experiments that supplied it
are reported in Appendix~\ref{sec:exp1}.
